\section{A repaired grouping construction and its geometric obligations}

This chapter replaces the arbitrary tile-to-square assignment in
\S\S3.1--3.4 of the paper.  It applies to simple polygonal disks, including
nonconvex polygons.  It does not assume convex prototiles.  Throughout,
collinear original vertices are retained on the boundaries of grouped tiles.
Consequently a grouped side remains a complete original side and scattering
does not introduce a side-to-side defect.

The construction uses one central group and $4N$ annular groups, instead of
the $4(N+1)$ independent small squares in the original drawing.  The four
corner squares are absorbed into adjacent annular groups.  Every new annular
group meets the centre along one connected arc and supplies one print.
There are still exactly $4N$ marked central sides, as required later.

All Lean names below are planned declaration names.  No declaration in this
chapter is claimed to have been checked by Lean.  The two formerly open
geometric estimates are proved below using polygonal Jordan theory,
piecewise-linear annulus routing, and an explicit compactness argument.
Those classical helper lemmas still require Lean implementations; no
geometric certificate is accepted as an assumption of the final theorem.

\begin{definition}[Mesh constants and the finite cellulation]
\label{def:mesh}
\lean{Tilings.Geometry.MeshConstants}
\uses{def:polygon,def:protoset,def:placed,def:tiling,def:sts}
Let $P$ be a finite family of full-dimensional simple polygonal disks.  Choose
$D>0$, $a>0$, $\gamma>0$, and $K_0\in\mathbb N$ such that every original
tile has diameter at most $D$, area at least $a$, every original side has
length at least $\gamma$, and every original tile has at most $K_0$ sides.

A finite STS patch covering a polygonal disk $\Omega$ is regarded as a finite
planar cellulation: original tiles are its closed faces, original sides are
its edges, and original vertices are its vertices.  A face interior is
connected, and distinct face interiors are disjoint.  We write
$B_\infty(r)=[-r,r]^2$ and use Minkowski sums $Q+B_\infty(r)$ for square
buffers.  A Euclidean diameter bound $D$ also bounds the coordinate diameter.
\end{definition}

The constants exist because $P$ is finite.  Neither FLC nor the crown-cover
lemma is used to obtain them.  In particular the grouping construction must
not depend on the final soficity theorem, on the desired routing certificate,
or on the FLC conclusion derived from that certificate.

\begin{lemma}[Regular grouping with an inside anchor]
\label{lem:regular-grouping}
\lean{Tilings.Geometry.regularGrouping}
\uses{def:mesh,def:tiling}
Let a finite STS patch of mesh $D$ cover $\Omega$, and let $Q$ be a square
with $Q+B_\infty(3D)\subset\operatorname{int}\Omega$.  There is a finite
union $L$ of complete original faces such that
\[
 Q\subset L\subset Q+B_\infty(D),
\]
$L$ is a polygonal disk, and the dual graph of its faces is connected.
Every edge of $\partial L$ adjoins on its inside an original face $F$ such
that $\operatorname{int}F\cap\operatorname{int}Q\ne\varnothing$.
\end{lemma}

\begin{proof}
Select the faces whose interiors meet $\operatorname{int}Q$ and let their
union be $U$.  The open square is contained in $\operatorname{int}U$.
Every selected face interior is connected and meets that square; thus
$\operatorname{int}U$ is connected.  The union is finite and regular closed.
Since every selected face has diameter at most $D$,
$U\subset Q+B_\infty(D)$.

Add every bounded connected component of the \emph{open} complement
$\mathbb R^2\setminus U$, taking its closure together with $U$.  A connected
face interior disjoint from $U$ is contained in one complementary component.
The added region therefore consists of complete original faces; no original
face is cut.  There are finitely many such components, since their boundaries
are contained in the finite original skeleton.

The complement outside the convex square $Q+B_\infty(D)$ is connected to
infinity, so every filled component lies inside that square.  This gives the
support bounds.  Filling a complementary component joins its face interiors
to the existing interior across boundary edges.  Hence the filled interior
is connected, and its open complement is the unbounded complementary
component and is connected.

Here is the classical planar fact needed in the formalization.  A finite
regular closed union of polygonal faces, with connected interior and
connected open complement, has a single simple polygonal boundary.  At a
boundary vertex there cannot be alternating blocks of inside and outside
faces: interior arcs joining the two inside blocks, together with a small
local arc, would separate the two outside blocks by the polygonal Jordan
theorem.  The outside could not be connected.  Therefore the boundary graph
has degree two at every boundary vertex.  It is a union of disjoint cycles;
connected interior and connected exterior exclude more than one cycle.
There are no zero-area appendages, since the region is a union of complete
faces.  Thus the boundary is one simple cycle and $L$ is a polygonal disk.

Filling deletes interior boundary edges and creates no new exterior edge.
The inside face along every remaining boundary edge was selected for $U$;
this proves the anchoring statement.  Finally join two face-interior points
by a generic polygonal path in $\operatorname{int}L$.  Perturb it away from
the finite vertices.  Its consecutive faces share complete sides, giving
a walk in the face-dual graph.
\end{proof}

For Lean, the proof should be divided into a finite-face complement lemma,
a polygonal Jordan-boundary lemma, and a generic-path-to-dual-walk lemma.
The Jordan-boundary lemma is a genuine topological result and must be proved
or obtained from a proved polygonal topology API; the definition of a
``group'' must not simply declare the filled union to be a disk.

\begin{lemma}[The dual outside-anchored cycle]
\label{lem:outside-grouping}
\lean{Tilings.Geometry.outsideAnchoredGrouping}
\uses{def:mesh,lem:regular-grouping}
Let $Q_1+B_\infty(10D)\subset\operatorname{int}\Omega$ and put
$K_1=Q_1+B_\infty(3D)$.  There is a polygonal disk $H$, made of complete
original faces, such that
\[
 Q_1+B_\infty(2D)\subset H\subset Q_1+B_\infty(3D).
\]
Every edge of $\partial H$ adjoins on its outside a face whose interior
meets the open exterior of $K_1$.
\end{lemma}

\begin{proof}
Let $E$ be the closed exterior of $\Omega$, together with every original
face whose interior meets the open exterior of $K_1$.  The open exterior
of $K_1$ is contained in $\operatorname{int}E$, and every selected face
interior meets it.  Thus $\operatorname{int}E$ is connected.

No selected face can enter $\operatorname{int}(Q_1+B_\infty(2D))$:
one of its coordinates would have to change by more than its diameter $D$
between that point and its point outside $K_1$.  Let $H$ be the closure of
the component of $\mathbb R^2\setminus E$ containing $Q_1$.  This gives
the lower inclusion.  The open exterior of $K_1$ lies in $E$, giving the
upper inclusion.  As in the preceding lemma, face interiors belong to
single complementary components, so $H$ is a union of complete faces.

Its interior is connected by its definition.  The open complement
$\mathbb R^2\setminus H$ is \emph{not} simply $\operatorname{int}E$: it
also contains the \emph{pockets}, the other components of
$\mathbb R^2\setminus E$, which consist of unselected faces.  Each pocket
borders selected faces across whole edges, so it is joined to
$\operatorname{int}E$ through those edges, and $\operatorname{int}E$ is
connected and disjoint from $H$.  Hence $\mathbb R^2\setminus H$ is
connected and $H$ has no holes.  A pocket face cannot sit across an edge
of $\partial H$, because the far side of such an edge lies in
$\operatorname{int}E$.  With connected interior and connected open
complement, alternating boundary sectors are excluded by the same
polygonal Jordan argument as before.  Hence $H$ is a polygonal disk.  Each exterior
boundary edge adjoins a face selected for $E$, proving the outside anchor.
The margin $10D$ ensures that this face is an original face of the patch
rather than the artificial unbounded exterior region.
\end{proof}

The distinction between inside and outside anchors matters.  A cycle lying
in a thin band around a square may run backwards along the band for a long
distance.  Its mere Hausdorff proximity to the square does not give a
uniform diameter bound for sectors between selected endpoints.

\begin{lemma}[Disjoint skeletal spokes]
\label{lem:annular-spokes}
\lean{Tilings.Geometry.existsAnnularSpokes}
\uses{def:mesh,lem:regular-grouping,lem:outside-grouping}
Let $N\ge2$ and $R$ be sufficiently large compared with $D$.  Take concentric
squares $Q_0,Q_1$ of side lengths $NR,(N+2)R$ respectively, with the margin
required above.  Construct $L,H$ by the preceding lemmas.  The annulus
$H\setminus\operatorname{int}L$ admits $4N$ mutually disjoint simple arcs
in the original skeleton, each joining $\partial L$ to $\partial H$ and
meeting those cycles only at its endpoints.  Their endpoints are distinct
and appear in the same cyclic order on the two cycles.  (The index $j$ of
the $j$th cut disk is matched with the ideal sector $T_j$ only in the next
lemma, through the walls: each wall crosses $\partial L$ exactly once, and
pigeonhole over the ideal wedges fixes the correspondence.)  The components
obtained by cutting along the arcs are $4N$ polygonal disks consisting of
complete original faces.  Each has a connected face-dual graph and shares
one connected nondegenerate arc with $\partial L$.
\end{lemma}

\begin{proof}
Mark the $4N$ positions subdividing the four sides of $Q_0$ into intervals
of length $R$, including the corners.  Draw the radial segments from the
square centre through these positions to the corresponding positions of
$K_1=Q_1+B_\infty(3D)$.  In the square annulus, successive radial segments
have separation at least a universal positive multiple of $R$.  This
follows by considering a straight side and the two sides adjacent to a
corner; the homothety factor is $1+2/N+6D/(NR)$.  Extending each segment
a bounded number of mesh diameters past the two cycles preserves disjoint
tubes when $R/D$ is sufficiently large.

Follow the finite sequence of original faces crossed by one extended
segment.  Consecutive faces share a side or vertex, and every face boundary
is a connected polygonal cycle.  Their boundary graphs form a connected
finite graph, all of whose edges lie within $D$ of the segment.  It meets
both $\partial L$ and $\partial H$.  Select a graph path, erase loops,
and take the part from its last point on $\partial L$ to its first
subsequent point on $\partial H$.  The part lies in the annulus and has
no other point on either cycle.  Planarity of the skeleton ensures that
an edge cannot cross a cycle without meeting it at a vertex.

The separated tubes give disjoint arcs and distinct endpoints.  Disjoint
arcs joining the two boundary components of an annulus have compatible
cyclic endpoint order.  Cutting a polygonal annulus along such $4N$ arcs
gives $4N$ disks, one between each successive pair; this is a direct
polygonal Jordan argument, or the elementary polygonal annulus theorem.
All cuts use original edges, so the disks are unions of whole original
faces.  Their inner boundary arcs have distinct endpoints and positive
length.  Their face-dual graphs are connected by the generic-path argument
used in Lemma~\ref{lem:regular-grouping}.
\end{proof}

\begin{lemma}[Quantitative sector localization and private neighbours]
\label{lem:quantitative-sector}
\lean{Tilings.Geometry.quantitativeAnnularGrouping}
\uses{def:mesh,lem:regular-grouping,lem:outside-grouping,lem:annular-spokes}
There are finite positive constants $C_{\rm sep},C_{\rm loc},C_{\rm len}$
and $C_{\rm err}$, independent of $N$, the substitution power, and the
particular STS patch, with the following property.  For
$R\ge C_{\rm sep}D$, the cycles and spokes above can be selected so that
each annular group $S_j$ lies in the $C_{\rm loc}D$ neighbourhood of the
corresponding ideal radial trapezoid $T_j$.  Along its outer boundary it
has at least
\[
 b(R,D)=\left\lfloor {R\over C_{\rm len}D}\right\rfloor-C_{\rm err}
\]
distinct fine edges whose outside neighbouring original faces meet no
other annular group $S_t$, $t\ne j$.  In particular, after choosing $R$
so that $b(R,D)\ge K_0+2$, every small group has at least $K_0+2$ ordinary
marked inner sides and at least three private neighbour occurrences
relative to every other small group.

There is consequently a constant $C$ independent of $N$ such that
$\operatorname{diam}S_j\le CR$, and
\[
 Q(P,R)=K_0\left\lceil {4C^2R^2\over a}\right\rceil
\]
bounds the number of retained sides of every small group.
\end{lemma}

\begin{proof}
Translate the common square centre to zero.  Put $A=NR/2$,
$h=A+R+3D$, and $\lambda=h/A$.  Let $v_j$ be the successive length-$R$
positions on $\partial Q_0$; the corresponding outer positions are
$\lambda v_j$ on $\partial K_1$.  Assume $R\ge1000D$.

An inner spoke endpoint $u_j$ satisfies
$A\le\|u_j\|_\infty\le A+D$ and lies within Euclidean distance $D$
of the ray through $v_j$.  If $|u_j-tv_j|\le D$, then
$A-D\le tA\le A+2D$.  Since $|v_j|\le\sqrt2A$,
\[
 |u_j-v_j|\le D+2\sqrt2D<4D.
\]
The outer endpoint $w_j$ satisfies
$h-D\le\|w_j\|_\infty\le h$, giving similarly
$|w_j-\lambda v_j|<4D$.

Use the inside-anchored face at $u_j$ to join it through that face's
interior to $p_j\in\operatorname{int}Q_0$.  The collar lies in the
$5D$ ball about $v_j$.  Join $p_j$ radially to zero inside $Q_0$.
Use the outside-anchored face at $w_j$ to reach
$q_j\notin K_1$, in the $5D$ ball about $\lambda v_j$; continue its
radial ray to infinity outside the convex square $K_1$.

These continuations are disjoint.  On a straight side, adjacent ideal
positions have determinant $AR$.  Perturbations of size at most $5D$
change it by at most $10\sqrt2AD+25D^2$, strictly less than $AR$.
The corner cases are rotations of the same estimate.  Applying it also
throughout the linear perturbation preserves cyclic angular order.
The endpoint boxes are separated by at least $R-10D$; the other radial
segments are separated there by at least $R/4$ before perturbation.
Thus a collar cannot meet another continuation.  Remove any self-loop
within its own endpoint box.  We obtain proper simple walls $W_j$ from
zero to infinity, pairwise disjoint except at zero.  In
$H\setminus\operatorname{int}L$ their portions are precisely the
skeletal spokes and stay within $5D$ of the ideal radial rays.

Let $T_j$ be the ideal trapezoid between successive radial segments and
the square boundaries $Q_0,K_1$.  For a point
$x\in\operatorname{int}K_1\setminus Q_0$ outside the two $5D$ radial
tubes, truncate $W_j,W_{j+1}$ on a sufficiently large circle and close
them along that circle.  The resulting polygonal Jordan loop determines
its wedge.  Deform the core portions inside $Q_0$, straighten the
annular portions inside their convex radial tubes, and deform the
exterior portions outside $K_1$.  The three pieces of each wall are glued
at the two points where the wall crosses $\partial L$ and $\partial H$;
the three homotopies are chosen to agree at these transition points, so
they assemble to one homotopy of the whole loop.  The exterior of $K_1$
is not simply connected, and the exterior deformation is legitimate only
because the exterior portion of a wall is a collar of diameter at most
$5D$ followed by a radial ray: it does not wind around $K_1$, so it is
homotopic, within the exterior of $K_1$ and through curves of increasing
infinity norm, to the ideal radial ray of the same angular order.  None
of these homotopies meets $x$.  Winding number is unchanged, so the actual and ideal wedge
labels agree away from the tubes.  This proves localization there.

If $x$ lies in one of the tubes, clip its nearest radial-line point to
the segment between infinity norms $A$ and $h$.  The clipping moves it
by at most $5\sqrt2D$.  Its distance to the corresponding boundary
segment of $T_j$ is therefore less than
$(1+\sqrt2)5D<13D$.  Taking closures gives the safe bound
\[
 S_j\subset T_j+B(0,20D).
\]
This argument does not assume that the square-band cycles have short or
monotone boundary arcs.

Take the middle half $J$ of the outer edge of $T_j$; its length is
$\lambda R/2\ge R/2$.  Its distance to every other ideal trapezoid is
at least $R/(4\sqrt2)$, by the distance to either bordering radial line.
The inward normal from each $y\in J$ meets $\partial H$ within its
first $D$ of length: $H$ contains the square with halfside $h-D$ and
is contained in the square with halfside $h$.  The middle interval is
at least $R/4$ from the outer corners.  Except at finitely many projected
vertices the intersection lies in an original fine edge $e$.

That intersection has distance at least
$R/(4\sqrt2)-21D>D$ from every other $S_t$.  Its outside original
face, of diameter at most $D$, cannot meet another small group.  Thus
$e$ is an ordinary marked private edge of $S_j$.  The projections of
these private edges cover $J$ apart from finitely many points, and each
has projected length at most $D$.  There are consequently at least
$|J|/D\ge R/(2D)$ distinct private edge occurrences.  Repeated outside
neighbour faces do not affect this occurrence count.

We may take $C_{\rm sep}=1000$, $C_{\rm loc}=20$,
$C_{\rm len}=2$, and $C_{\rm err}=0$.  Since
$\lambda\le2+3D/R$, each ideal trapezoid has span at most $5R$ along
its square side and at most $2R$ across it (the roles of the coordinates
swap for the side sectors), hence diameter less than $6R$.
Localization gives $\operatorname{diam}S_j\le7R$.  A square of side
$14R$ contains $S_j$, so disjoint constituent interiors of area at
least $a$ bound their number by $\lceil196R^2/a\rceil$.
Thus $Q(P,R)=K_0\lceil196R^2/a\rceil$ is admissible.  Retaining all
original vertices proves the stated side bound.
\end{proof}

\begin{lemma}[Inner sector order]
\label{lem:inner-sector-order}
\lean{Tilings.Geometry.quantitativeAnnularGrouping_innerOrder}
\uses{def:mesh,lem:regular-grouping,lem:quantitative-sector}
Let $N\ge2$, $R\ge1000D$, and let $L$ be the inside-anchored group of $Q_0$
with $Q_0\subseteq L\subseteq Q_0+B_\infty(D)$, listed counterclockwise
(winding number $2\pi$ about every interior point).  Let
$y_1,\dots,y_m\in\partial L$ with $y_k$ in the $20D$-neighbourhood of the
ideal sector $T_{j_k}$, where $j_1<\dots<j_m$, and with $y_k$ at distance at
least $50D$ from every other ideal sector $T_t$, $t\ne j_k$.  Then
$y_1,\dots,y_m$ appear in this cyclic order along the counterclockwise
boundary parametrization of $L$.
\end{lemma}

\begin{proof}
This is the inside-anchored counterpart of the outer order statement in the
proof of Lemma~\ref{lem:quantitative-sector} (points of $\partial H$ near
the extended radial segments appear in index order along $\partial H$), and
the remark before Lemma~\ref{lem:annular-spokes} explains why anchoring is
needed: a cycle in the band $Q_0+B_\infty(D)\setminus\operatorname{int}Q_0$
may run backwards along the band for a long distance, so Hausdorff proximity
to $Q_0$ alone does not order its points.  Join every $y_k$ through its
inside-anchoring face (a face of $L$ of diameter at most $D$ whose interior
meets $\operatorname{int}Q_0$) to a point $z_k\in\operatorname{int}Q_0$ with
$|z_k-y_k|\le D$, and then radially to the centre $c$.  These paths lie in
$L$, stay within $D$ of the sectors $T_{j_k}$, which are at least $50D$
apart at the points $y_k$, and meet only at $c$; they cut the disk $L$ into
$m$ sectors whose boundary arcs between consecutive points $y_k,y_{k+1}$
contain no other $y_i$.  A boundary parametrization visiting the points out
of sector order would have to run through the perimeter parameter (the
radial projection to $\partial Q_0$) backwards by at least one sector gap,
which the winding separators of the outer order argument, transplanted to
the inner band with the radial escape paths above in place of the
outside-anchored escape faces, exclude; a parametrization running clockwise
contradicts the winding number $2\pi$ of the counterclockwise listing.  The
constants are those of the Lean statement: $20D$ localization, $50D$
separation, $R\ge1000D$.
\end{proof}

Choose $R$ using that lemma before choosing $N$.  Then choose $N\ge2$ with
$4N>3Q(P,R)$.  Only afterwards choose $k$ so that every $\tau^k$-macrotile
contains the required squares and margin.  Each original polygon has a
chosen interior ball of positive radius; scaling that ball suffices, even
for a nonconvex polygon.  There is no circular dependence of $Q$ on $N$.

Apply the grouping independently to the finitely many fixed patches
$\tau^k\alpha$, $\alpha\in P$.  Their finitely many resulting labelled
groups form $P'$.  This argument needs no claim that arbitrary patches have
only finitely many geometric neighbourhoods.  Such a claim is needed later
for FLC clones and must follow separately from
Lemma~\ref{lem:crown-cover-supertile}, whose proof is independent of routing.

\begin{definition}[Finite geometric macrocertificate]
\label{def:certificate}
\lean{Tilings.Geometry.MacroCertificate}
\uses{def:polygon,def:protoset,def:placed,def:substitution,def:sts,def:mesh}
A macrocertificate consists of finite form and type sets, a finite
side-to-side polygonal template for each form, and the following explicit
finite data and properties.
\begin{enumerate}
\item[(G1)] Every template covers the scaled parent disk, has connected
full inner-side dual graph, and has at least three children.  Every type
has an inner side.  Type tags of different templates are disjoint.
Each internal template edge has its own unique interface identifier,
carried by both incident owners; an inner side is carried by exactly two
children of its template (the side and its partner occurrence).  It records the template, the two type
tags, their incident declared-side slots, and their relative placement;
different interfaces between the same two types remain distinguishable.
In the decorated implementation the identity is
stored in the first component except on a print, where it is stored in
the fifth component.  A special side which is not a print still uses the
first component; specialness alone does not move identity to the fifth.
\item[(G2)] The boundary belt graph is connected.  Its vertices are
\emph{all} types having a vertex on the macroboundary, including types
touching it only at a vertex.  Its edges are inner sides having an endpoint
on the macroboundary.  Each boundary-belt type has such an edge.  The
central child has neither outer nor boundary sides.
\item[(G3)] For each parent declared side, take the ordered list of whole
fine outer sides contained strictly inside its scaled side interval.
At most two fine sides are omitted for crossing its endpoints.  There
are five distinct ports in this list, chosen by a fixed interior-quantile
convention.  Ranks and port numbers are counted in one global positive
direction fixed for each unoriented supporting-line direction, never in
the boundary traversal order of the owner: the quantiles
$\lceil iL/6\rceil$ are not symmetric under reversal (for $L=12$ they are
$2,4,6,8,10$ one way and $11,9,7,5,3$ the other), and the properness
lemma needs port $i$ of a side to face port $i$ of the coincident
opposing side.  Since only translations act, both owners of a coincident
side compute the same direction.  Full matching of a parent declared
interval and its fine tiling then induces identical contained lists and
port numbers.  A parent
dummy endpoint may lie inside a fine child side: no nested subdivision
or hereditary-cut assumption is made.  The finite template also records
the interval intersections with each parent declared side and with the
declared boundary partitions at all required intermediate levels.
Their exact endpoint/frame encoding and matching implications are proved
in the boundary-transport chapter; scalar ordinal ranks alone do not
establish them for arbitrary collinear dummy vertices.
\item[(G4)] Marked parent sides have distinct central prints in cyclic
order.  Cyclic order is the index order of the listed sides, which agrees
with the planar order only when all listed polygons are traversed in the
same rotational sense; the certificate therefore lists \emph{every} form
of its alphabet, the original singleton forms included, counterclockwise
(an original form whose given listing is clockwise is relisted; this
changes side indices but no carrier, side set or scattering).  For every marked side $d$ there are four paths $P_i(d)$ from port
$i$ to its print.  Every side occurrence is used by at most one path,
except the terminal print in either of its two incident owners.  Thus
the central and small owners of that same print are both exempt at the
shared terminal side.  Every path uses marked sides only.  \emph{A path is a sequence of geometric sides, so a genuine path contains both occurrences (both owners) of every inner side it crosses}, the terminal print excepted: the path records the print in its small owner only.  Without this clause the third cell $A_i(d)$ would sit on one occurrence of a crossed side and $0$ on its partner, and the lifting theorem fails for the same reason as below (found by Aristotle, 12 October 2026).  \emph{Every
side of a genuine path other than its initial port is an inner side of the
template}: a path never passes through an outer side of another macroside.
Without this clause the lifting theorem fails, because a path crossing an
outer side $b$ of an adjacent macroside would put a nonzero parent signal
into the third cell of $b$ while the coincident side of the neighbouring
block, lying on no path, carries zero.  The only central side on a genuine
path is its final print.  In its small owner the print is the unique
special side of that owner (the one side of a small type whose cells are
copied from the parent; this is the exception allowed in (G6)).  Its preceding
side is an ordinary side of the small child sharing that print.  The
fifth path and all five paths of unmarked parent sides are singleton
ports, disjoint from genuine paths.  This is a finite routing certificate.
\item[(G5)] There is a unique large central child and all other forms are
original or small.  The central-form map fixes every retained large form.
For a large parent the print map is the identity on its marked sides.
There are $L=4N$ marked central sides, at most $Q$ marked sides on a
noncentral form, and $L>3Q$.  Every original type has a marked inner side.
\item[(G6)] Noncentral fingerprints distinguish types.  For an original
type its fingerprint is the multiset of type components of its ordinary
inner-side identity pairs.  For a small type use its ordinary marked-side
identity pairs and allow at most two further type components from its
unique special side, which is an inner side (the side shared with the
central child).  The possible fingerprint multisets for two
distinct noncentral types are disjoint.  This is a finite separating
property, not an assertion that a composed decorated tile already exists.
\item[(G7)] Marking discipline.  Every outer side, every boundary side
(inner side with an endpoint on the macroboundary), every print, and every
side occurring on a genuine path is marked.  Consequently an unmarked side
is an inner, nonprint, nonboundary side lying on no path; in the
construction these are exactly the sides shared by the central child and
its ring other than the prints.  This condition is what makes the
``local cells'' of Definition~\ref{def:control/typed-decoration} genuinely
type-determined: without it an unmarked boundary side, whose fourth cell
is the parent header, or an unmarked path side, whose third cell is a
parent signal, would carry a parent-dependent value declared local, which
is the error type of the source article.  The paper's own instantiation
satisfies (G7), since it marks every side not belonging to a large
quasi-square.
\end{enumerate}
\end{definition}

The certificate does \emph{not} assume the final soficity statement,
existence of a coloured composition, or a matching-rule system with the
desired language.  These are the conclusions proved from the finite data
in the combinatorial chapter.  Conversely, the final geometric theorem
must construct the certificate from the hypotheses on $\tau$; accepting
the certificate as an unproved input would only prove a conditional core.

\begin{lemma}[A finite sufficient fingerprint criterion]
\label{lem:fingerprint}
\lean{Tilings.Geometry.fingerprintSeparation}
\uses{def:certificate,lem:quantitative-sector}
Suppose (G1) holds, every original form has at most $K_0$ sides, every small
type has at least $K_0+2$ ordinary marked inner sides and one special side,
and, for each pair of distinct small types in one template, the first has
at least three ordinary marked-side occurrences adjacent to original
types which are not ordinary marked-side neighbours of the second.
Then (G6) holds.
\end{lemma}

\begin{proof}
An original fingerprint has at most $2K_0$ components.  A small fingerprint
has at least $2(K_0+2)$ components before its special contribution.  Thus
an original and a small type cannot have equal fingerprints.

Ordinary components of types in different templates lie in disjoint tag
sets.  Every original fingerprint is nonempty.  A small type has more
than two ordinary components, whereas another template can supply at most
two such foreign components through its special side.  This separates
types in different templates.

Consider two original types $u,v$ in one template.  In the fingerprint of
$u$, the multiplicity of $u$ is its inner-side degree, and the multiplicity
of $v$ is the number of its sides adjacent to $v$.  Equal fingerprints
force both degrees to equal the number of $u$--$v$ sides.  All their inner
sides therefore join them to one another.  They form a two-vertex connected
component, contradicting (G1).

For two distinct small types in one template, count the three private
original-neighbour occurrences of the first.  None occurs in the ordinary
fingerprint of the second.  Its special side can add at most two type
components, so it cannot supply all three.  Repeated occurrences of the
same private neighbour are allowed: the argument is about multiset
multiplicity, not distinct neighbour names.  The fingerprints differ.
\end{proof}

\begin{lemma}[Wide routing after a sufficiently large power]
\label{lem:routing}
\lean{Tilings.Geometry.existsRoutingAfterPower}
\uses{def:mesh,def:certificate,lem:quantitative-sector,lem:finite-idempotent}
For the finite grouped family $P'$ constructed above, extend the grouped
substitution $\sigma$ to every group by scattering its original
constituents and applying $\sigma$ to them.  Choose central children as
in the paper and let $c:P'\to P'$ be their form map.  There is an
arbitrarily large power exponent $l$ with $c^{2l}=c^l$ such that, after
retaining the chosen central child and its annular small neighbours and
scattering all other groups in each $\sigma^l$-template, the finite data
required by (G4) exist; the print of every marked parent side is one of the
$4N$ designated marked sides of the central form (the side it shares with one
annular piece), distinct prints on distinct pieces, in cyclic and planar
matching order with the marked parent sides, and the identity on a large
parent.  The setup also records the sector localization of the annular
pieces (Lemma~\ref{lem:quantitative-sector}) and the piece of every print;
the pieces of the prints of the marked sides are pairwise distinct and in
the cyclic order of the marked sides (Lemma~\ref{lem:inner-sector-order},
applied in the construction of Lemma~\ref{lem:certificate-exists}).  The
input is the annular structure of the retained
templates from Lemma~\ref{lem:quantitative-sector}: the central child is the
inside-anchored group $L$, the small neighbours are the $4N$ annular groups
between $L$ and the outside-anchored group $H$ (so every private outer edge
adjoins a face reaching the exterior of $K_1$), $R\ge10000(K_0+3)D$, and
every annular group has four private outer edges in separated windows,
pairwise at distance at least $R/10$; all grouped disks are listed
counterclockwise (as in the construction of
Lemma~\ref{lem:certificate-exists}), so that the polygon loop of a parent
is the positively oriented boundary parametrization used in the
matching-order clauses.  Without the anchoring, the private
faces of a small neighbour could lie in cavities closed off by fewer than
four sides, and the four disjoint paths cannot exist.
\end{lemma}

\begin{proof}
The finite idempotence lemma permits arbitrarily large suitable $l$.
The final chosen centre lies inside the scaled first-generation central
child.  That first-generation child is strictly inside its parent
template.  The finitely many parent templates have a positive minimum
clearance.  After normalizing a $\sigma^l$-template by its expansion
$s=\theta^{kl}$, its central positions consequently lie in a fixed compact
set $K_A\subset\operatorname{int}A$, uniformly over $l$.
For the finite parent catalogue put
\[
 \rho=\min_A\operatorname{dist}(K_A,\partial A)>0.
\]
This handles the moving centre without assuming convergence or convexity.

Outside the retained central annulus every child is an original tile,
so the mesh remains $D$ and every fine outer side has length in
$[\gamma,D]$.  A scaled parent declared side of length $s\ell$ contains
an ordered list of whole fine sides of cardinality
$L\ge s\ell/D-2$: only the at most two sides crossing its endpoints
are omitted.  Choose five ports by the sixth-quantiles
$\lceil iL/6\rceil$, $i=1,\ldots,5$, for sufficiently large $L$.
Their separation from one another and from the ends is bounded below by
\[
 \gamma\bigl(\lfloor L/6\rfloor-1\bigr)-D.
\]
After normalization this gives a positive lower bound depending only on
the finite parent catalogue.  Closed portions of different sides away
from their endpoints also have positive mutual Euclidean separation:
they are disjoint compact subsets of a simple polygonal boundary.
Thus all realized genuine port tuples belong to a compact set of ordered
boundary configurations with uniform separation, including separation
from every other designated or singleton port.  No dummy endpoint is
assumed to be a fine child vertex.

In each selected small group's middle outer interval, choose four
separated windows, for example at relative parameters
$0.30,0.43,0.57,0.70$, and use the projection argument of
Lemma~\ref{lem:quantitative-sector} to choose a private fine edge in
each window.  Let $p_i$ be its midpoint.  Its outside-anchored original
face contains a point $q_i$ outside $K_1$; join $p_i$ to $q_i$ through
that face.  This collar has diameter at most $D$.  The windows are
separated by a fixed multiple of $R$, so the four faces are distinct.
Taking $R\ge10000(K_0+3)D$ makes all these collars and radial rays from
$q_i$ to infinity separated by more than $10D$ outside $K_1$.
The rays remain outside $H$.  These choices are fixed for the finitely
many central annulus catalogues.  Reverse the four entrance labels
where the port numbering is opposite to counterclockwise boundary order;
the paired lists then have the same cyclic order.

For a normalized centre $z\in K_A$, take the intermediate square
$C_z=z+[-\rho/4,\rho/4]^2$.  Its circumradius is less than $\rho/2$,
so it lies strictly in $A$.  Each fixed radial direction meets its
boundary at a point $b_i(z)$.  Consider the polygonal annulus
$A\setminus\operatorname{int}C_z$.  Its ordered outer port points and
inner points $b_i(z)$ admit disjoint simple polygonal arcs.  To prove this
classical helper, cut the annulus along a polygonal crosscut, apply the
piecewise-linear disk theorem, and route ordered lines in the resulting
strip.  Explicitly: for nested simple closed polygons $C\subset\operatorname{int}A$
there is a continuous map $\Psi\colon[0,1]\times\mathbb R\to A\setminus\operatorname{int}C$,
$1$-periodic in the second coordinate and injective on $[0,1]\times[0,1)$,
sending $\{0\}\times\mathbb R$ onto $\partial A$ and $\{1\}\times\mathbb R$ onto
$\partial C$ as loop parametrizations and $(0,1)\times\mathbb R$ onto
$\operatorname{int}A\setminus C$.  It is obtained from the polygonal disk
theorem (ear decomposition of a simple polygon, gluing of disk models along
a diagonal, Alexander-trick boundary reparametrization) applied to the
annulus cut open along one crosscut, by identifying the two copies of the
crosscut; the strip $[0,1]\times[0,1]$ is the cut-annulus model.  The
entrance points and directions below are not a finite catalogue fixed in
advance: the routing statement is uniform over \emph{all} admissible
entrance data, so the compactification argument must include, besides the
normalized centre and the whole port segments, the radial entrance
directions, and the cyclic order of the entrance pieces must be derived
from the planar order of the prints on the contact arcs (the
\texttt{print\_order} datum), not from their indices.  Concretely, the
setup records the piece of every print, and these pieces are in the cyclic
order of the marked sides (Lemma~\ref{lem:inner-sector-order}); every
entrance point lies within $21D$ of the ideal sector of its piece (the
sector localization of Lemma~\ref{lem:quantitative-sector} plus one face of
diameter at most $D$) and outside $K_1$, and entrance points of distinct
marked sides are at least $R/10$ apart, so the radial directions of the
entrances of distinct marked sides are in the cyclic order of the sides
(entrances of adjacent sectors cannot cross their common ray: two points
outside $K_1$ within $21D$ of that ray and of $\partial K_1$ are within
$100D<R/10$ of each other); the four entrances of one marked side are
relabelled so that their radial order matches the port order.  Short disjoint inward-normal outer collars and outward-radial
inner collars can be prescribed first: remove thin attached fingers
around them, route between their tips in the remaining annulus, and
concatenate.  This does not require $A$ to be star shaped.

For one parameter tuple the compact arcs have positive mutual distance
and positive distance from the annulus boundary off these endpoint
collars.  We now prove local stability of these bounds.  A small movement
of a port along its open side is realized in a local rectangle by a PL
slide $(t,n)\mapsto(t+\delta\psi(t,n),n)$, with $\psi=1$ on its
collar and zero at the rectangle boundary.  If
$|\delta|\operatorname{Lip}\psi<1/2$, this is bi-Lipschitz and
preserves the parent disk and the straight normal collar.  The supports
for different ports are disjoint and avoid the central region.

A small change $z\mapsto z'$ is realized by
\[
 H(x)=x+\phi(x)(z'-z),
\]
where a PL bump $\phi$ is one on $B(z,\rho/2)$ and zero outside
$B(z,3\rho/4)$.  The support lies in $\operatorname{int}A$.  Fix the
normalized collar length $\kappa<\rho/4$ once and for all; then the slide
rectangles at the outer ports and the bump support $B(z,3\rho/4)$ are
disjoint from each other and from the outer collars, so the bump fixes
every outer collar and each slide fixes the inner square.
For $|z'-z|\operatorname{Lip}\phi<1/2$, the lower Lipschitz bound
gives injectivity and the contraction equation
$x=y-\phi(x)(z'-z)$ gives surjectivity.  Thus $H$ is a bi-Lipschitz
homeomorphism, translates $C_z$ and its short radial collars exactly,
and fixes $\partial A$.  Compose these maps and shrink the parameter
neighbourhood so the total distortion is at most two.  The old routes
then give routes for every nearby tuple with at least half the old
clearance and the required straight collars.

The centre and port parameter space is compact and there are only
finitely many parent, annulus, and pairing patterns.  Routes are built,
and the local stability argument is applied, for every parameter in the
\emph{closure} of the set of realized tuples, not only for realized ones;
this is legitimate because the separation bounds defining the parameter
space are closed conditions, and it is needed because a finite subcover
is taken of that compact closure.  The subcover therefore gives a uniform
$c_0>0$ for mutual and off-collar boundary clearance in normalized
coordinates.  Choose $l$ larger, still with
$c^{2l}=c^l$, so that $c_0s>100D$ and the physical intermediate square
contains the fixed central annulus.  Join each inner radial ray at its
intermediate-square endpoint to the outer PL route, whose initial
inner collar has the same radial direction.  The two portions meet
only at that endpoint, since one is inside the square and the other
outside it.  Different radial portions are separated by a multiple of
$R$; different outer portions by a multiple of $s$.  Mixed pairs are
also separated: middle outer portions stay $c_0s$ from the whole square,
and distinct endpoint collars are separated from the other rays by the
finite direction separation scaled by $s$.

The first $O(D)$ part of an outer normal collar crosses only faces in
an $O(D)$ ball about its own port, which contains no other port tile.
Afterwards it is more than $D$ from the parent boundary.  The collar may
cross several ordinary original faces; it need not remain in its own
port face.  It avoids every other designated and singleton port tile,
which is the condition actually required.

Finally perturb the routes away from the finite original vertices and
list their crossed faces.  Consecutive faces share whole sides.  Erase
repeated-face loops while preserving the first port face and final
outside-anchored face.  Begin the resulting side sequence with the
port, end it with the private entry edge, and append the print.  The
entry and print belong to the same small group, hence are consecutive
sides of a path in the paper's exact sense.  Different routes cannot
cross the same original face: their $D$ tubes are disjoint away from
the collars, their inner collar faces are distinct, and their outer
collar balls are separated.  Thus every nonterminal side occurrence
is used at most once.  The shared terminal print is exempt in both
its central and small owners.  No central side is met earlier.
All traversed original sides, small entry sides, and prints are marked.
Every singleton port path remains disjoint from the genuine routes.
\end{proof}

\begin{lemma}[Construction of the geometric macrocertificate]
\label{lem:certificate-exists}
\lean{Tilings.Geometry.existsMacroCertificate}
\uses{def:substitution,def:sts,def:certificate,lem:regular-grouping,lem:outside-grouping,lem:annular-spokes,lem:quantitative-sector,lem:fingerprint,lem:routing,lem:finite-idempotent,lem:inner-sector-order}
For every STS polygonal substitution $\tau$, there are a finite grouped
family $P'$, a scattering map $\Delta$, an integer $m>0$, and a grouped
substitution $\sigma'$ satisfying
\[
 \Delta\sigma'\alpha=\tau^m\Delta\alpha
 \quad\hbox{for all }\alpha\in P',
\]
together with a macrocertificate (G1)--(G7).  The grouped family contains
every original form as a singleton group and every group is a group of tiles
of a supertile (the hypotheses of
Lemma~\ref{lem:grouping-lift}), and when $\tau$ has hereditary listed
sides so has $\sigma'$.
\end{lemma}

\begin{proof}
Choose $R$, then $Q$, then $N$, then $k$ in the noncircular order described
above.  Group each fixed $\tau^k\alpha$ by the inside/outside cycles and
the spokes.  Keep the finite labelled groups in $P'$.  Their supports are
disks and their constituent dual graphs are connected.  Scattering a group
and applying the grouped substitution to each constituent defines $\sigma$
and gives $\Delta\sigma=\tau^k\Delta$.

\emph{Step 3 (the routing setup).}  For every grouped form $\alpha$ the
template $\sigma\alpha$ is the regrouping of $\tau^k\Delta\alpha$: each
constituent $q\in\Delta\alpha$ contributes the translated groups of its
supertile $\tau^kq$ (its $L_q$, its $4N$ annular pieces, singletons).  The
\emph{central child} $c(\alpha)$ of $\sigma\alpha$ is the group $L_q$ of the
constituent $q$ whose centre $z_q$ is nearest to the centre of $\alpha$; thus
$c$ maps grouped forms to large catalogue forms $L_p$.  For a power $l$ the
retained template of $\alpha$ keeps, inside $\sigma^l\alpha$, only the
central chain $c^l(\alpha)$ (the central child of the central child, $l$
times) and its $4N$ annular neighbours, and scatters every other group into
original tiles.  Its annulus is a translate of the catalogue annulus of one
original supertile, so the quantitative data ($N$, $R$, the four separated
windows of every piece, the private-edge count) are those of the grouping
step.  The clearance of the central child grows geometrically along the
central chain: at each step the chosen centre lies in the square of
Lemma~\ref{lem:quantitative-sector} with its margin, strictly inside the
scaled parent, so after normalization by $\theta^{kl}$ the central child
stays at a positive distance $\rho$ from the boundary of the retained region,
uniformly in $l$ and $\alpha$ (a geometric series with ratio $\theta^{-k}$).
The marked sides are fixed by the marking rule: every side of a form that is
not a central image is marked; a large form $L_p$ has exactly $4N$ marked
sides, the side it shares with each of the $4N$ annular pieces (one per
piece; these are whole sides of $L_p$ by the whole-side property of groups).
The print slots of $\alpha$ send its marked sides injectively to contact
sides of the central child, preserving cyclic order and in planar matching
order with the marked sides (both disks are listed counterclockwise, and an
order-preserving injection between two cyclically ordered boundary families
is in matching cyclic order in the annulus between them); this is possible
because a non-large form has at most $Q<4N/3$ sides while the centre has $4N$
contact sides, and for a large parent $\alpha=c^l\beta$ the print map is the
identity (its $4N$ marked sides are its own contact sides).  The prints
are chosen sector by sector: for the marked sides $a_1<\dots<a_m$ in index
order take distinct pieces $j_1<\dots<j_m$ (all $4N$ pieces for a large
parent) and for each $j_k$ the contact side of the central child met by the
ray from the centre through the middle of the sector $T_{j_k}$; by the
sector localization the face across that side lies in the piece $j_k$, so
the piece map preserves cyclic order by construction, and by
Lemma~\ref{lem:inner-sector-order} these sides occur counterclockwise in
the order of the sectors, so the print map preserves cyclic order as well;
the sector localization itself is retained from
Lemma~\ref{lem:quantitative-sector}.

\emph{Step 3$'$ (the routed power).}  Lemma~\ref{lem:routing} applied to this
setup gives $l>0$ with $c^{2l}=c^l$ (Lemma~\ref{lem:finite-idempotent}) and
routing data for every retained template at the power $l$.  Keep only central
large forms in $\operatorname{im}c^l$.  For a retained large form
$\alpha=c^l\beta$, $c^l\alpha=c^{2l}\beta=c^l\beta=\alpha$.  Thus the central
form is fixed.  Put $m=kl$ and call the resulting substitution $\sigma'$.
Scattering commutes with its decomposition, proving the displayed
compatibility.

\emph{Step 4 (the finite certificate data).}  The forms of the certificate
are the retained catalogue tiles (the central chain forms, the annular
pieces, and the original singletons), each with its disk relisted
counterclockwise and translated to the origin.  The types are the positioned
children of the retained templates; the owner of a type is its template and
its anchor is its position.  A type is large, small or original according as
its form is a central image, an annular piece or an original singleton.  The
partner of an inner side occurrence is the other occurrence of the same
geometric side in the same template (templates are side-to-side, so an inner
side has exactly two owners, Step 5); the interface identifier of an inner
side is its unordered pair of occurrences, with the two type tags recorded in
a fixed order.  Marking follows the marking rule of Step 3 through the parent
sides.  The special side of a small type is its designated contact side with
the central child: the print when the type receives one (prints land on
distinct pieces), an arbitrary fixed contact side otherwise.  The print of a
marked parent side is the routing's print slot.  The outer list of a parent
declared side is the list of whole outer fine sides of the template contained
in the scaled side, without repetition, in the lexicographically positive
direction of the supporting line.  The genuine paths are the occurrence
chains of the routed side lists: for a routed side list $e_0$ (the port),
$e_1,\dots,e_r$ (the print) with crossed children $X_0,\dots,X_{r-1}$ the
chain is $(X_0,e_0),(X_0,e_1),(X_1,e_1),(X_1,e_2),\dots,(X_{r-1},e_r)$: both
occurrences of every crossed inner side, the print in its small owner only;
its side projection with consecutive repetitions removed is the routed list.
The fifth path and all paths of unmarked sides are singleton ports.  The
chains inherit from the routing data that every side after the port is
inner, that every side is marked, that no central side occurs before the
print, that the side before the print is an ordinary side of the small
owner, that no occurrence repeats, and that distinct paths share only the
terminal print in either owner.  The central child of a template is owned
by that template; the mixed form of a type's child is the catalogue form of
its own form; each mixed form has one catalogue entry; and the relabelling
between the counterclockwise listing of a certificate disk and the listing
of its grouped disk is a bijection of side indices.  The
constants are $N$ and $Q$ as chosen, and the rank bound is the largest outer
list length.  This data represents $\tau$ at the power $m$: $\Delta\sigma'=
\tau^{kl}\Delta$ because scattering commutes with the retained/scattered rule;
the dissections are side-to-side with connected duals, whole sides and
retained vertices because they are catalogue groups; every original form is a
singleton; every group lies in a supertile; heredity transfers.  The retained
templates are side-to-side because scattering a group of a side-to-side patch
along whole sides preserves the side-to-side property.

\emph{Step 5 (G1).}  Disk cellulations with at least the centre and its ring
have connected inner-side duals and at least three children.  Every child has
an inner side: a child whose sides are all outer is the whole template.  An
inner side has exactly two owners, because the template is side-to-side and
three polygons with disjoint interiors cannot share a whole segment.  The
partner and identifier clauses are the definitions of Step 4.

\emph{Step 6 (G2).}  Around the macroboundary, consecutive boundary faces are
connected through the finite fan at each boundary vertex.  Its connecting
edges are precisely inner sides with an endpoint on that boundary.  Faces
touching the boundary only at a vertex must be included in this fan.  Every
boundary type has such an inner edge because it is not the whole macrotile.
The retained centre is strictly interior, by the clearance of Step 3.

\emph{Step 7 (G3).}  The lexicographic order is a global positive
orientation.  The finite boundary interval intersections give the outer
lists; at most two fine sides are omitted for crossing the endpoints of a
scaled side (Lemma~\ref{lem:whole-ports}); each list has at least five
entries because the five ports of the routing are distinct entries of it
(they are the quantiles of this very list); the lists are bounded by the rank
bound.  Ranks and ports are numbered in the global supporting-line direction,
as (G3) requires.  No nested-cut assumption is used.  The exact matching
implications of the endpoint/frame schedules are proved separately in the
boundary-transport chapter.

\emph{Step 8 (G4).}  The print map is the routing's print slot map: injective
and cyclic-order preserving on the marked sides, with inner prints whose other
owner is a small child, where the print is the special side.  The occurrence
chains of Step 4 inherit from the routing that every side after the port is
inner, that every side is marked (sides of non-large children and prints),
that no central side occurs before the print, that the side before the print
is an ordinary side of the small owner (the entry edge and the print belong
to the same small group), that no occurrence repeats, and that distinct paths
share only the terminal print in either owner (their crossed faces are
distinct).

\emph{Step 9 (G5).}  The central child belongs to its template and is its
unique large child.  For a retained large form $\alpha=c^l\beta$,
$c^l\alpha=\alpha$, so the central form is fixed and the print map is the
identity on its marked sides (the central child has the parent's own disk).
A large form has exactly $4N$ marked sides (Step 3); a non-large form has at
most $Q$ marked sides, because all its sides are marked and an annular piece
has at most $Q$ sides (Lemma~\ref{lem:quantitative-sector}) while an original
tile has at most $K_0\le Q$ sides; $4N>3Q$ by the choice of $N$.  Every
original type has a marked inner side: all its sides are marked and it has an
inner side (Step 5).

\emph{Step 10 (G6).}  The special sides are those of Step 4: inner, shared
with the centre, unique per small type.  Every small type has at least
$K_0+2$ private outer edges (Lemma~\ref{lem:quantitative-sector}), so the
private-neighbour estimate and Lemma~\ref{lem:fingerprint} give the
disjointness of the possible fingerprint multisets of distinct noncentral
types.

\emph{Step 11 (G7).}  The marking rule ``every side not belonging to the large
central group, plus the prints'' gives (G7): outer and boundary sides belong
to non-large children (the centre is strictly interior), prints are marked,
and every occurrence on a genuine path lies in a non-large child, since no
central side is met before the print.

This constructs the finite geometric certificate.  Neither a coloured
composition theorem nor the final soficity conclusion has been assumed.
All Lean implementations, including the classical polygonal topology
helpers used in this chapter, remain unverified.
\end{proof}

\paragraph{Why a citation to an earlier matching-rule theorem is insufficient.}
Fernique--Ollinger's theorem assumes a \emph{good} combinatorial
substitution; its connecting condition is the very information-routing
condition at issue here.  The implication from arbitrary planar STS data
to that condition must still be proved.  Goodman--Strauss assumes
hereditary sides and sibling-side-to-side conditions, and directly treats
substitution tilings.  Using it as an alternate route requires verification
of those exact hypotheses and a proved theorem supporting the hierarchical
variant as well.  Neither a bibliographic citation nor an unproved Lean
axiom supplies these constructions or their Lean implementations.

\paragraph{Triangulation is not an unproved refinement reduction.}
The suggested polygonal triangulation is only an auxiliary tool in the
continuous wide-routing proof.  Replacing prototiles by triangles would
require a separate finite substitution-compatible refinement and a proof
of projection equivalence for both substitution and hierarchical tilings.
No such reduction is assumed here.
